% ============ 25 Lefschetz 不动点公式 (原书 pp.147-149) ============
\chapter{Lefschetz 不动点公式}\label{sec:25}
\index{Lefschetz 不动点公式}
\index{迹}
\index{对应}
\index{横截相交}
\index{正规相交}

这里我们证明：Lefschetz 不动点公式的存在性是具有良好性质的环映射，K\"unneth 公式与 Poincar\'e 对偶存在的形式推论.本节中，$X$ 是代数闭域 $k$ 上的非奇异簇(通常完备).

设 $V$ 是向量空间，$\varphi \colon V \to V$ 是线性映射.若 $(a_{ij})$ 是 $\varphi$ 关于基 $(e_i)$ 的矩阵，则 $\varphi$ 的迹(记为 $\operatorname{Tr}(\varphi|_{V})$)是 $\sum a_{ii}$.它不依赖于基的选取.若 $(f_i)$ 是对偶向量空间 $\check{V}$ 的对偶基，使得 $e_i \cdot f_j = \delta_{ij}$，则
\[
\sum_{i} \varphi(e_i) \cdot f_i = \sum_{i} \left( \sum_{j} a_{ji} e_j \right) \cdot f_i = \sum a_{ii} = \operatorname{Tr}(\varphi|_{V}).
\]

\begin{theorem}\label{thm:25.1}[Lefschetz 不动点公式]
设 $X$ 是代数闭域 $k$ 上的完备非奇异簇，$\varphi \colon X \to X$ 是正则映射.则
\[
(\Gamma_{\varphi} \cdot \Delta) = \sum (-1)^{r} \operatorname{Tr}(\varphi^{*}|H^{r}(X, \mathbb{Q}_{\ell})).
\]
这里 $\Gamma_{\varphi}$ 是 $\varphi$ 的图象，$\Delta$ 是 $X \times X$ 中的对角.于是 $(\Gamma_{\varphi} \cdot \Delta)$ 是按重数计的 $\varphi$ 的不动点数.
\end{theorem}

\begin{example}\label{ex:25.2}
考虑映射
\[
\varphi \colon \mathbb{P}^{1} \to \mathbb{P}^{1}, \qquad (x_0 \colon x_1) \mapsto (x_0 + x_1 \colon x_1).
\]
在 $x_1 \neq 0$ 的仿射片上是 $x \mapsto x + 1$，在 $x_0 \neq 0$ 的仿射片上是 $x \mapsto \frac{x}{1+x}$.由于 $\varphi$ 在 $H^{0}(\mathbb{P}^{1})$ 与 $H^{2}(\mathbb{P}^{1})$ 上都恒等作用(因为它次数为 1), Lefschetz 不动点公式的右边是 2.$\varphi$ 的唯一不动点是 $1 = (1 \colon 0)$($x_0 \neq 0$ 的仿射片中的原点).为计算 $(\Gamma_{\varphi} \cdot \Delta)_1$，必须计算曲线
\[
y(1+x) = x, \qquad y = x
\]
在 $(0, 0)$ 处的交数.它是
\[
k[x, y]/(y - x, y(1+x) - x) = k[x]/(x + x^{2} - x) = k[x]/(x^{2})
\]
的维数，即 2，正如定理~\ref{thm:25.1} 预测的.
\end{example}

为简化叙述，我们固定同构 $\mathbb{Q}_{\ell} \to \mathbb{Q}_{\ell}(1)$；这相当于对所有 $n$ 选取相容的同构 $\mathbb{Z}/\ell^{n}\mathbb{Z} \to \mu_{\ell^{n}}(k)$.

写 $H^{*}(X) = \bigoplus H^{r}(X, \mathbb{Q}_{\ell})$---它是 $\mathbb{Q}_{\ell}$-代数.K\"unneth 公式允许我们把 $H^{*}(X \times X)$ 等同于 $H^{*}(X) \otimes H^{*}(X)$，即把 $p^{*}(a) \smile q^{*}(b)$ 等同于 $a \otimes b$.这里 $p$ 与 $q$ 是投影映射 $X \times X \rightrightarrows X$.

Poincar\'e 对偶给出非退化配对
\[
H^{*}(X) \times H^{*}(X) \to H^{2d}(X, \mathbb{Q}_{\ell}) \cong \mathbb{Q}_{\ell}.
\]
写 $e_{2d}$ 为 $H^{2d}(X)$ 的典范生成元(任何点 $P$ 的类).

下面定理~\ref{thm:25.1} 的证明几乎一字不差地从一本拓扑书上抄来.

设 $\varphi \colon X \to Y$ 是正则映射.下一个引理表明映射 $\varphi^{*} \colon H^{*}(Y) \to H^{*}(X)$ 等于 $X \times Y$ 上对应 $\Gamma_{\varphi}$ 定义的映射.

\begin{lemma}\label{lem:25.3}
对任何正则映射 $\varphi \colon X \to Y$ 与任何 $y \in H^{*}(Y)$，
\[
p_{*}(\mathrm{cl}_{X \times Y}(\Gamma_{\varphi}) \smile q^{*}y) = \varphi^{*}(y).
\]
\end{lemma}

\begin{proof}
我们计算：
\[
p_{*}(\mathrm{cl}_{X \times Y}(\Gamma_{\varphi}) \smile q^{*}y) = p_{*}((1, \varphi)_{*}(1) \smile q^{*}y) \quad \text{(24.2b)}
\]
\[
= p_{*}(1, \varphi)_{*}(1 \smile (1, \varphi)^{*}q^{*}y) \quad \text{(24.2e)}
\]
\[
= (p \circ (1, \varphi))_{*}(1 \smile (q \circ (1, \varphi))^{*}y) \quad \text{(24.2c)}
\]
\[
= \mathrm{id}_{*}(1_X \smile \varphi^{*}y)
\]
\[
= \varphi^{*}(y).
\]
\end{proof}

\begin{lemma}\label{lem:25.4}
设 $(e_i)$ 是 $H^{*}(X)$ 的一组基，$(f_i)$ 是 $H^{*}(X)$ 中相对于杯积对偶的基，使得 $e_i \smile f_j = \delta_{ij} e_{2d}$($\delta_{ij}$ 是 Kronecker 记号).对任何正则映射 $\varphi \colon X \to X$,
\[
\mathrm{cl}_{X \times X}(\Gamma_{\varphi}) = \sum \varphi^{*}(e_i) \otimes f_i.
\]
\end{lemma}

\begin{proof}
由于 $f_i$ 作为 $\mathbb{Q}_{\ell}$-向量空间构成 $H^{*}(X)$ 的基，它们也作为 $H^{*}(X)$-模构成 $H^{*}(X \times X) = H^{*}(X) \otimes_{\mathbb{Q}_{\ell}} H^{*}(X)$ 的基.因此，
\[
\mathrm{cl}_{X \times X}(\Gamma_{\varphi}) = \sum a_i \otimes f_i
\]
对唯一元素 $a_i \in H^{*}(X)$ 成立.由引理~\ref{lem:25.3},
\[
\varphi^{*}(e_j) = p_{*}\left( \left( \sum_{i} a_i \otimes f_i \right) \smile (1 \otimes e_j) \right) = p_{*}(a_j \otimes e_{2d}) = a_j.
\]
\end{proof}

\begin{proof}
(定理的证明) 设 $e^{r}_i$ 是 $H^{r}$ 的一组基，$f^{2d-r}_i$ 是 $H^{2d-r}$ 的对偶基.则
\[
\mathrm{cl}(\Gamma_{\varphi}) = \sum_{r, i} \varphi^{*}(e^{r}_i) \otimes f^{2d-r}_i, \qquad \text{且}
\]
\[
\mathrm{cl}(\Delta) = \sum_{r, i} e^{r}_i \otimes f^{2d-r}_i = \sum_{r, i} (-1)^{r(2d-r)} f^{2d-r}_i \otimes e^{r}_i = \sum_{r, i} (-1)^{r} f^{2d-r}_i \otimes e^{r}_i.
\]
对这两个表达式取乘积，发现
\[
\mathrm{cl}_{X \times X}(\Gamma_{\varphi} \cdot \Delta) = \sum_{r, i} (-1)^{r} \varphi^{*}(e^{r}_i) f^{2d-r}_i \otimes e_{2d} = \sum_{r} (-1)^{r}\operatorname{Tr}(\varphi^{*}|H^{r})(e_{2d} \otimes e_{2d}).
\]
现在对两边应用 $\varphi_{X \times X}$.
\end{proof}

\begin{remark}\label{rem:25.5}
尽管在上面的讨论中我们等同了 $\mathbb{Q}_{\ell}$ 与 $\mathbb{Q}_{\ell}(1)$，上述定理无需这个等同也照原样成立.要点是
\[
H^{r}(X, \mathbb{Q}_{\ell}(s)) = H^{r}(X, \mathbb{Q}_{\ell}) \otimes \mathbb{Q}_{\ell}(s),
\]
且 $\varphi^{*}$ 穿过 $H^{r}(X, \mathbb{Q}_{\ell})$ 作用.与一维 $\mathbb{Q}_{\ell}$-向量空间 $\mathbb{Q}_{\ell}(s)$ 作张量积不改变迹.
\end{remark}

对 $\varphi$ 的不动点 $P$，有一个判别 $(\Gamma_{\varphi} \cdot \Delta)_P = 1$ 的准则是有用的.

设 $Y$ 与 $Z$ 是非奇异簇 $X$ 的闭子簇，并假定点 $P$ 是 $Y \cap Z$ 的不可约分支.则当
\begin{enumerate}
  \item $Y$ 与 $Z$ 在 $P$ 处非奇异，
  \item $\operatorname{Tgt}_{P}(Y) \cap \operatorname{Tgt}_{P}(Z) = 0$，且
  \item $\dim Y + \dim Z = \dim X$
\end{enumerate}
时，$(Y \cdot Z)_P = 1$.条件 (b) 意味着 $Y$ 与 $Z$ 在 $P$ 处横截相交，条件 (c) 意味着 $Y$ 与 $Z$ 在 $P$ 处恰当相交(即 $\operatorname{codim} P = \operatorname{codim} Y + \operatorname{codim} Z$).

\begin{lemma}\label{lem:25.6}
设 $\varphi \colon X \to X$ 是正则映射，$P \in X$ 是 $\varphi$ 的不动点.则当 $1$ 不是 $(d\varphi)_P \colon \operatorname{Tgt}_{P}(X) \to \operatorname{Tgt}_{P}(X)$ 的特征值时，$(\Gamma_{\varphi} \cdot \Delta)_P = 1$.
\end{lemma}

\begin{proof}
把前面的评注应用于 $\Gamma_{\varphi} \cap \Delta$ 上的点 $(P, P)$.因为 $\Gamma_{\varphi}$ 与 $\Delta$ 都同构于 $X$，条件 (a) 与 (c) 成立.因为 $\operatorname{Tgt}_{(P,P)}(\Gamma_{\varphi})$ 是 $(d\varphi)_P \colon \operatorname{Tgt}_{P}(X) \to \operatorname{Tgt}_{P}(X)$ 的图象，而 $\operatorname{Tgt}_{(P,P)}(\Delta)$ 是恒等映射 $\operatorname{Tgt}_{P}(X) \to \operatorname{Tgt}_{P}(X)$ 的图象，条件 (b) 成立，当且仅当 $1$ 不是 $(d\varphi)_P$ 的特征值.
\end{proof}
